% Korean edition of figs/layers.tex.
\fig[wide]{layers}{REST 요청은 HTTP/2 연결 하나를 공유하고, 실시간 피드는 두 번째 WebSocket 연결이 필요하다}[채널별 클라이언트 측 프로토콜 스택: 애플리케이션부터 연결 설정까지]{%
\begin{tikzpicture}[x=1mm,y=1mm,
  box/.style={draw=navy,line width=0.5pt,fill=white,minimum width=30mm,minimum height=8mm,inner sep=0pt,font=\fontsize{7.5}{9}\selectfont,anchor=north west},
  lname/.style={anchor=west,align=left,font=\fontsize{7.5}{9}\selectfont\bfseries\color{navy}},
  note/.style={anchor=west,align=left,text width=69mm,inner sep=0pt,font=\fontsize{7}{8.4}\selectfont\color{greyD}},
  hd/.style={anchor=south,align=center,font=\fontsize{7.5}{9}\selectfont\bfseries\color{navy}}]
% column heads
\node[hd] at (58,-1.2) {REST\\[-0.5mm]{\fontsize{7}{8.4}\selectfont\mdseries\color{greyD}\code{/api}}};
\node[hd] at (92,-1.2) {실시간 피드\\[-0.5mm]{\fontsize{7}{8.4}\selectfont\mdseries\color{greyD}\code{/hubs/main}}};
% stacks: rows of 8 mm from y = -2
\foreach \y/\n/\a/\b in {-2/애플리케이션/{REST API}/SignalR, -10/프레이밍/{HTTP/2 스트림}/WebSocket, -18/HTTP/{HTTP/2}/{HTTP/1.1}, -26/TLS/{TLS 1.3}/{TLS 1.3}, -34/TCP/TCP/TCP}{
  \node[box] at (43,\y) {\a};
  \node[box] at (77,\y) {\b};
  \node[lname] at (23,\y-4) {\n};
}
% DNS, before connecting
\node[box,draw=greyD,dashed,minimum width=64mm] at (43,-45) {DNS};
\node[lname,text=greyD] at (23,-49) {연결 전};
% bracket over TLS, TCP and DNS of both stacks
\draw[navy,line width=0.6pt] (21.5,-26.3) -- (20,-26.3) -- (20,-52.7) -- (21.5,-52.7);
\draw[navy,line width=0.6pt] (20,-39.5) -- (18.5,-39.5);
\node[anchor=east,align=right,font=\fontsize{7}{8.4}\selectfont\bfseries\color{navy}] at (18,-39.5) {새 연결마다\\다시 치르는\\비용};
% short notes, one per row
\node[note] at (111,-6)  {REST: 요청마다 Bearer 토큰.\\피드: 연결을 열 때 토큰을 한 번 검사};
\node[note] at (111,-14) {REST: 요청마다 스트림 하나, 연결 하나에 여러 개.\\피드: JSON 레코드, 각각 0x1E로 끝남};
\node[note] at (111,-22) {REST: ALPN으로 선택. 피드: 협상용 HTTP/2 POST 후,\\별도 연결에서 \code{Upgrade: websocket}};
\node[note] at (111,-30) {왕복 1회. 피드는 자체 핸드셰이크를 따로 함};
\node[note] at (111,-38) {REST: 15\textasciitilde 30\,s 유휴 후 서버가 닫음.\\피드: 3\,s마다 요청해 열어 둠};
\node[note] at (111,-49) {\code{contest.timefolio.net} $\rightarrow$ 221.168.38.40};
\end{tikzpicture}
}{상자 하나가 프로토콜 하나다. 두 채널 모두 같은 DNS 조회에서 시작한다. 이후 REST 요청은 같은 연결에 새 HTTP/2 스트림을 열어 DNS, TCP, TLS를 건너뛴다. 서버의 HTTP/2 SETTINGS 프레임에 \code{ENABLE_CONNECT_PROTOCOL}이 없어 브라우저를 포함한 어떤 클라이언트도 HTTP/2 위의 WebSocket[1]을 쓸 수 없으므로, 실시간 피드에는 두 번째 연결이 필요하다. 서버에서는 nginx 프런트엔드가 요청을 ASP.NET 애플리케이션으로 넘긴다. SETTINGS 프레임은 Python \code{h2} 라이브러리로 읽었다. 프런트엔드는 오류 페이지로, 애플리케이션은 오류 메시지로 식별했다.}{DART Quant Team}
